% !PRINTED 529
\weberpage{529}
%<<CONT
as a function belonging to $Q_1$, then the two values conjugate to it are
\weq{7}{\begin{aligned}
y_1 &= (\alpha-\alpha_2)(\alpha_3-\alpha_1)\\
y_2 &= (\alpha-\alpha_3)(\alpha_1-\alpha_2).
\end{aligned}}
These three quantities are the roots of a cubic partial resolvent, whose coefficients are easily formed as invariants of $f$ (cf.\ \S.\,64, 65, where the quantities $y$, $y_1$, $y_2$ were denoted by $U$, $V$, $W$). The equation becomes
\weq{8}{y^3 - Ay + \sqrt{D} = 0.}
Its discriminant is $B^2$, and for the product of the root differences the sign can easily be determined (by comparison of the terms of highest order according to the \gesperrt{Gauss} definition, \S.\,45, or also from \S.\,64). Thus one obtains:
\weq{9}{(y-y_1)(y-y_2)(y_1-y_2) = B.}
The equation (8) is, as already remarked, a normal equation, and it is easy to represent $y_1$, $y_2$ actually in terms of $y$. If one multiplies (9) by $(y_1-y_2)$, one finds, by making use of the formulae $(y-y_1)(y-y_2) = 3y^2 - A$, $(y_1-y_2)^2 = y^2 - 4y_1y_2 = 4A - 3y^2$:
\weq{10}{
\begin{aligned}
B(y_1-y_2) &= -9y^4 + 15Ay^2 - 4A^2\\
 &= 6Ay^2 + 9y\sqrt{D} - 4A^2\\
y_1+y_2 &= -y,
\end{aligned}
}
whereby $y_1$, $y_2$ are determined. If therefore one further adjoins $y$, the group of the biquadratic equation reduces to
\[
Q_1 = 1, \ (0,2)(1,3), \ (0,1)(2,3), \ (0,3)(1,2).
\]
This group is imprimitive (and indeed in three ways), and the equation can now be solved by two square roots. This is best reached in the following way. The three quantities
\weq{11}{
\begin{aligned}
v_1 &= (\alpha+\alpha_1-\alpha_2-\alpha_3)^2\\
v_2 &= (\alpha-\alpha_1+\alpha_2-\alpha_3)^2\\
v_3 &= (\alpha-\alpha_1-\alpha_2+\alpha_3)^2
\end{aligned}
}
all admit the permutations of the group $Q_1$ and are therefore rationally representable in terms of $y$. One obtains this representation as follows: We have
%CONT>>
